\documentclass[12pt, a4paper]{article}

% ─── Paquetes ────────────────────────────────────────────────────────────────
\usepackage[utf8]{inputenc}
\usepackage[T1]{fontenc}
\usepackage[american]{babel}
\usepackage{amsmath, amssymb, amsthm}
\usepackage{mathtools}
\usepackage{geometry}
\usepackage{graphicx}
\usepackage{xcolor}
\usepackage{booktabs}
\usepackage{array}
\usepackage{fancyhdr}
\usepackage{titlesec}
\usepackage{enumitem}
\usepackage{mdframed}
\usepackage{tcolorbox}
\usepackage{hyperref}


\tcbuselibrary{theorems, skins, breakable}

% ─── Geometría ───────────────────────────────────────────────────────────────
\geometry{
  a4paper,
  top=2.5cm,
  bottom=2.5cm,
  left=3cm,
  right=3cm
}

% ─── Colores institucionales ─────────────────────────────────────────────────
\definecolor{azulprimario}{RGB}{0, 71, 133}
\definecolor{azulclaro}{RGB}{210, 228, 245}
\definecolor{grisclaro}{RGB}{245, 245, 245}
\definecolor{verdedestacado}{RGB}{0, 120, 80}
\definecolor{rojoalerta}{RGB}{180, 30, 30}

% ─── Estilo de página ────────────────────────────────────────────────────────
\pagestyle{fancy}
\fancyhf{}
\fancyhead[L]{\small\color{azulprimario}\textit{Taller 02 — Difusión de Calor en un Disco}}
\fancyhead[R]{\small\color{azulprimario}\textit{Ecuaciones Diferenciales Parciales}}
\fancyfoot[C]{\small\thepage}
\renewcommand{\headrulewidth}{0.4pt}
\renewcommand{\headrule}{\hbox to\headwidth{\color{azulprimario}\leaders\hrule height \headrulewidth\hfill}}

% ─── Títulos de sección ──────────────────────────────────────────────────────
\titleformat{\section}[block]
  {\large\bfseries\color{azulprimario}}
  {\thesection.}{0.5em}{}
  [\vspace{0.5ex}\color{azulprimario}\hrule\vspace{0.5ex}]

\titleformat{\subsection}[block]
  {\normalsize\bfseries\color{azulprimario!80!black}}
  {\thesubsection.}{0.5em}{}

% ─── Cajas de entorno ────────────────────────────────────────────────────────
\tcbset{
  resultadoestilo/.style={
    enhanced,
    colback=azulclaro,
    colframe=azulprimario,
    fonttitle=\bfseries,
    title={Resultado},
    arc=3mm,
    boxrule=0.8pt,
    breakable
  },
  destacadoestilo/.style={
    enhanced,
    colback=grisclaro,
    colframe=verdedestacado,
    fonttitle=\bfseries\color{verdedestacado},
    arc=2mm,
    boxrule=0.6pt,
    breakable
  },
  notaestilo/.style={
    enhanced,
    colback=yellow!10,
    colframe=orange!70!black,
    fonttitle=\bfseries\color{orange!70!black},
    title={Nota},
    arc=2mm,
    boxrule=0.6pt,
    breakable
  }
}

\newenvironment{resultado}[1][]{%
  \begin{tcolorbox}[resultadoestilo, title={#1}]%
}{%
  \end{tcolorbox}%
}

\newenvironment{destacado}{%
  \begin{tcolorbox}[destacadoestilo]%
}{%
  \end{tcolorbox}%
}

\newenvironment{nota}{%
  \begin{tcolorbox}[notaestilo]%
}{%
  \end{tcolorbox}%
}

% ─── Operadores ──────────────────────────────────────────────────────────────
\newcommand{\dd}{\,\mathrm{d}}
\newcommand{\pd}[2]{\dfrac{\partial #1}{\partial #2}}
\newcommand{\pdd}[2]{\dfrac{\partial^2 #1}{\partial #2^2}}

% ─── Hyperref ────────────────────────────────────────────────────────────────
\hypersetup{
  colorlinks=true,
  linkcolor=azulprimario,
  urlcolor=azulprimario,
  pdftitle={Taller 02 - Difusión de Calor en un Disco},
  pdfauthor={Ecuaciones Diferenciales Parciales}
}

% ═══════════════════════════════════════════════════════════════════════════════
\begin{document}

% ─── PORTADA ─────────────────────────────────────────────────────────────────
\begin{titlepage}
  \centering
  \vspace*{1cm}

  {\color{azulprimario}\rule{\linewidth}{1.5pt}}\\[0.4cm]
  {\color{azulprimario}\rule{\linewidth}{0.4pt}}\\[1.5cm]

  {\huge\bfseries\color{azulprimario}
    Taller 02\\[0.4cm]
    \large Difusión de Calor en un Disco\\
    con Frontera Oscilante
  }\\[1.5cm]

  {\color{azulprimario}\rule{0.6\linewidth}{0.4pt}}\\[1cm]

  \begin{tcolorbox}[
    enhanced,
    colback=azulclaro,
    colframe=azulprimario,
    arc=4mm,
    boxrule=1pt,
    width=0.85\textwidth
  ]
  \centering
  {\large\bfseries Enunciado del Problema}\\[0.5em]
  {\normalsize
  Determinar el perfil de temperatura $u(r,t)$ en una placa circular de radio $R$,
  gobernado por la ecuación de calor en coordenadas polares con simetría radial,
  sujeto a una frontera oscilante $u(R,t) = T_0\sin(\omega t)$
  y temperatura inicial nula $u(r,0) = 0$.
  }
  \end{tcolorbox}

  \vspace{1.5cm}

  \begin{tabular}{rl}
    \textbf{\color{azulprimario}Curso:}    & Ecuaciones Diferenciales Parciales \\[0.3em]
    \textbf{\color{azulprimario}Taller:}   & 02 \\[0.3em]
    \textbf{\color{azulprimario}Método:}   & Expansión en Autofunciones de Fourier-Bessel \\[0.3em]
  \end{tabular}

  \vfill

  {\color{azulprimario}\rule{\linewidth}{0.4pt}}\\[0.4cm]
  {\color{azulprimario}\rule{\linewidth}{1.5pt}}
\end{titlepage}

% ─── TABLA DE CONTENIDOS ─────────────────────────────────────────────────────
\tableofcontents
\newpage

% ═══════════════════════════════════════════════════════════════════════════════
\section{Formulación del Problema}
% ═══════════════════════════════════════════════════════════════════════════════

Considérese una placa circular metálica delgada de radio $R$. El perfil de temperatura
$u(r,t)$ en la placa, bajo la hipótesis de simetría radial (independencia del ángulo $\theta$),
está gobernado por la \textbf{ecuación de calor en coordenadas polares}:

\begin{equation}
  \pd{u}{t} = \alpha^2 \left( \pdd{u}{r} + \frac{1}{r}\pd{u}{r} \right),
  \qquad 0 \leq r < R,\quad t > 0
  \label{eq:calor_original}
\end{equation}

donde $\alpha^2$ es la difusividad térmica del material. El problema queda completamente
determinado por las siguientes condiciones:

\begin{align}
  \textbf{(CF)}\quad &u(R,t) = T_0 \sin(\omega t) \label{eq:CF}\\[4pt]
  \textbf{(CI)}\quad &u(r,0) = 0 \label{eq:CI}
\end{align}

La condición de frontera \eqref{eq:CF} modela una \textit{fluctuación armónica} de temperatura
en el borde exterior del disco, con amplitud $T_0$ y frecuencia angular $\omega$.


% ═══════════════════════════════════════════════════════════════════════════════
\section{Paso 1: Transformación de Variable}
% ═══════════════════════════════════════════════════════════════════════════════

\subsection{Motivación}

El principal obstáculo para aplicar el método de expansión en autofunciones es que la condición
de frontera \eqref{eq:CF} es \textbf{no homogénea}. Los métodos espectrales requieren que en
la frontera la solución sea nula. Se introduce por tanto una transformación que \emph{traslade}
esta dificultad de la frontera hacia el interior de la ecuación.

\subsection{Definición de la nueva variable}

Se observa que la función $w(t) = T_0\sin(\omega t)$ satisface trivialmente la condición de
frontera en $r = R$. Se define entonces la nueva variable:

\begin{equation}
  \boxed{v(r,t) \;=\; u(r,t) - T_0\sin(\omega t)}
  \label{eq:transformacion}
\end{equation}

de donde se despeja:
\begin{equation}
  u(r,t) = v(r,t) + T_0\sin(\omega t)
  \label{eq:u_en_v}
\end{equation}

\subsection{Condiciones para la nueva variable}

\textbf{Condición de frontera:} evaluando \eqref{eq:transformacion} en $r = R$:
\begin{equation}
  v(R,t) = \underbrace{u(R,t)}_{T_0\sin(\omega t)} - T_0\sin(\omega t) = 0
\end{equation}

\textbf{Condición inicial:} evaluando \eqref{eq:transformacion} en $t = 0$:
\begin{equation}
  v(r,0) = \underbrace{u(r,0)}_{0} - T_0\sin(0) = 0
\end{equation}

\begin{resultado}[Nuevo problema — condiciones]
\begin{align*}
  v(R,t) &= 0 \qquad \text{(frontera homogénea)} \\
  v(r,0) &= 0 \qquad \text{(condición inicial nula)}
\end{align*}
\end{resultado}

\begin{nota}
La transformación \eqref{eq:transformacion} no elimina la no-homogeneidad: la mueve desde la
frontera hacia la EDP, donde aparecerá como un término fuente $G(t)$ (véase el Paso~2).
Este intercambio es conveniente porque una EDP no homogénea con frontera homogénea
admite solución en serie.
\end{nota}


% ═══════════════════════════════════════════════════════════════════════════════
\section{Paso 2: Nueva EDP para \texorpdfstring{$v(r,t)$}{v(r,t)}}
% ═══════════════════════════════════════════════════════════════════════════════

\subsection{Cálculo de las derivadas de \texorpdfstring{$u$}{u}}

Se sustituye \eqref{eq:u_en_v} en la EDP original \eqref{eq:calor_original}. Para ello
se calcula cada derivada por separado.

\medskip
\noindent\textbf{Derivada temporal:}
\begin{equation}
  \pd{u}{t} = \pd{}{t}\bigl[v(r,t) + T_0\sin(\omega t)\bigr]
            = \pd{v}{t} + T_0\,\omega\cos(\omega t)
\end{equation}

\noindent\textbf{Derivadas espaciales:} dado que $T_0\sin(\omega t)$ no depende de $r$,
sus derivadas respecto a $r$ son nulas:
\begin{align}
  \pdd{u}{r} &= \pdd{v}{r} \\[4pt]
  \frac{1}{r}\pd{u}{r} &= \frac{1}{r}\pd{v}{r}
\end{align}

Por tanto, el operador laplaciano en coordenadas polares satisface:
\begin{equation}
  \nabla^2 u \;=\; \pdd{u}{r} + \frac{1}{r}\pd{u}{r} \;=\; \nabla^2 v
\end{equation}

\subsection{Sustitución en la EDP original}

Reemplazando en \eqref{eq:calor_original}:
\begin{equation}
  \pd{v}{t} + T_0\,\omega\cos(\omega t) = \alpha^2\,\nabla^2 v
\end{equation}

Despejando $\pd{v}{t}$:
\begin{equation}
  \pd{v}{t} = \alpha^2\,\nabla^2 v - T_0\,\omega\cos(\omega t)
\end{equation}

\subsection{Identificación del término fuente}

Comparando con la forma canónica requerida $v_t = \alpha^2\nabla^2 v + G(t)$, se identifica:

\begin{resultado}[Término fuente]
\begin{equation}
  G(t) = -T_0\,\omega\cos(\omega t)
  \label{eq:Gt}
\end{equation}
\end{resultado}

\subsection{Resumen del problema transformado}

\begin{tcolorbox}[
  enhanced,
  colback=grisclaro,
  colframe=azulprimario,
  title={\textbf{Problema para $v(r,t)$}},
  fonttitle=\bfseries,
  arc=3mm, boxrule=0.8pt, breakable
]
\begin{align}
  \pd{v}{t} &= \alpha^2\,\nabla^2 v - T_0\,\omega\cos(\omega t),
  \quad 0 \leq r < R,\; t > 0 \label{eq:EDP_v}\\[4pt]
  v(R,t) &= 0 \label{eq:CF_v}\\[4pt]
  v(r,0) &= 0 \label{eq:CI_v}
\end{align}
\end{tcolorbox}


% ═══════════════════════════════════════════════════════════════════════════════
\section{Paso 3: Expansión en Series de Fourier-Bessel}
% ═══════════════════════════════════════════════════════════════════════════════

\subsection{Autofunciones del problema}

El problema espacial asociado a \eqref{eq:EDP_v}--\eqref{eq:CF_v} posee la misma estructura
que la ecuación de Bessel de orden cero. Las \textbf{autofunciones} que satisfacen la condición
de nulidad en la frontera son:
\begin{equation}
  \phi_n(r) = J_0\!\left(\frac{\lambda_n\, r}{R}\right), \qquad n = 1, 2, 3, \ldots
\end{equation}

donde $\lambda_n$ es el $n$-ésimo cero positivo de la función de Bessel de primera especie y
orden cero, definido por la condición:
\begin{equation}
  J_0(\lambda_n) = 0
\end{equation}

Los primeros valores son aproximadamente: $\lambda_1 \approx 2.4048$, $\lambda_2 \approx 5.5201$,
$\lambda_3 \approx 8.6537$.

\begin{nota}
Esta construcción es completamente análoga al caso cartesiano: en un dominio $[0,L]$
con frontera nula, las autofunciones son $\sin(n\pi r/L)$ y se anulan en $r = L$
porque $n\pi$ son los ceros del seno. Aquí, $\lambda_n$ son los ceros de $J_0$
y cumplen el mismo papel.
\end{nota}

\subsection{Propuesta de expansión}

Se propone escribir $v(r,t)$ como combinación lineal infinita de las autofunciones:
\begin{equation}
  \boxed{v(r,t) = \sum_{n=1}^{\infty} T_n(t)\, J_0\!\left(\frac{\lambda_n\, r}{R}\right)}
  \label{eq:expansion}
\end{equation}

donde:
\begin{itemize}[leftmargin=2em]
  \item $J_0\!\left(\tfrac{\lambda_n r}{R}\right)$ captura la \textbf{dependencia espacial} y está completamente determinada;
  \item $T_n(t)$ son los \textbf{coeficientes temporales} desconocidos que deben determinarse.
\end{itemize}

\subsection{Propiedad clave del Laplaciano sobre las autofunciones}

Cada autofunción $\phi_n$ satisface la ecuación de Bessel de orden cero, lo que implica:
\begin{equation}
  \nabla^2 J_0\!\left(\frac{\lambda_n\, r}{R}\right)
  = -\left(\frac{\lambda_n}{R}\right)^2 J_0\!\left(\frac{\lambda_n\, r}{R}\right)
  \label{eq:eigenvalor}
\end{equation}

Esto expresa que las autofunciones son \emph{eigenvalores} del operador laplaciano,
con valor propio $-(\lambda_n/R)^2$.

\subsection{Sustitución de la expansión en la EDP}

\noindent\textbf{Lado izquierdo} de \eqref{eq:EDP_v}:
\begin{equation}
  \pd{v}{t} = \sum_{n=1}^{\infty} T_n'(t)\, J_0\!\left(\frac{\lambda_n\, r}{R}\right)
\end{equation}

\noindent\textbf{Lado derecho} de \eqref{eq:EDP_v}, usando \eqref{eq:eigenvalor}:
\begin{equation}
  \alpha^2\,\nabla^2 v + G(t)
  = \sum_{n=1}^{\infty} T_n(t)\left[-\alpha^2\frac{\lambda_n^2}{R^2}\right]
    J_0\!\left(\frac{\lambda_n\, r}{R}\right) + G(t)
\end{equation}

El término fuente $G(t)$ también se expande en la misma base:
\begin{equation}
  G(t) = \sum_{n=1}^{\infty} g_n(t)\, J_0\!\left(\frac{\lambda_n\, r}{R}\right)
  \label{eq:expansion_G}
\end{equation}

donde $g_n(t)$ son los coeficientes de proyección de $G(t)$ sobre las autofunciones
(se calcularán en el Paso~4).

\subsection{Ecuación para cada coeficiente temporal}

Igualando ambos lados y usando la \textbf{ortogonalidad} de las autofunciones
(si dos series en la misma base son iguales, sus coeficientes son iguales término a término),
se obtiene una EDO lineal de primer orden para cada $T_n(t)$:

\begin{equation}
  T_n'(t) + k_n\, T_n(t) = g_n(t), \qquad n = 1, 2, 3, \ldots
  \label{eq:ODE_Tn}
\end{equation}

donde se define el parámetro:
\begin{equation}
  k_n = \alpha^2\left(\frac{\lambda_n}{R}\right)^2
  \label{eq:kn}
\end{equation}

\begin{resultado}[Reducción a EDOs]
La EDP bidimensional \eqref{eq:EDP_v} ha sido reducida a una familia de EDOs
lineales de primer orden \eqref{eq:ODE_Tn}, una para cada modo $n$.
La solución general de \eqref{eq:ODE_Tn} mediante factor integrante es:
\begin{equation}
  T_n(t) = e^{-k_n t}\int e^{k_n t}\, g_n(t)\, \dd t + C\,e^{-k_n t}
  \label{eq:solucion_general_Tn}
\end{equation}
La constante $C$ se determina con la condición inicial $T_n(0) = 0$.
\end{resultado}


% ═══════════════════════════════════════════════════════════════════════════════
\section{Paso 4: Coeficientes, Solución de la EDO y Ensamblaje Final}
% ═══════════════════════════════════════════════════════════════════════════════

\subsection{Parte A: Cálculo de \texorpdfstring{$g_n(t)$}{gn(t)} por ortogonalidad}

Para determinar el coeficiente $g_n(t)$ de la expansión \eqref{eq:expansion_G},
se multiplica ambos miembros de dicha expansión por $J_0\!\left(\tfrac{\lambda_m r}{R}\right) \cdot r$
y se integra de $0$ a $R$:

\begin{equation}
  \int_0^R G(t)\, J_0\!\left(\frac{\lambda_m r}{R}\right) r\,\dd r
  = \sum_{n=1}^{\infty} g_n(t) \int_0^R J_0\!\left(\frac{\lambda_n r}{R}\right)
    J_0\!\left(\frac{\lambda_m r}{R}\right) r\,\dd r
\end{equation}

Aplicando la \textbf{relación de ortogonalidad} del apéndice:
\begin{equation}
  \int_0^R J_0\!\left(\frac{\lambda_n r}{R}\right)
  J_0\!\left(\frac{\lambda_m r}{R}\right) r\,\dd r
  =
  \begin{cases}
    0 & \text{si } n \neq m \\[4pt]
    \dfrac{R^2}{2}\left[J_1(\lambda_n)\right]^2 & \text{si } n = m
  \end{cases}
\end{equation}

la suma infinita colapsa al único término $n = m$:
\begin{equation}
  \int_0^R G(t)\, J_0\!\left(\frac{\lambda_n r}{R}\right) r\,\dd r
  = g_n(t)\cdot \frac{R^2}{2}\left[J_1(\lambda_n)\right]^2
\end{equation}

Despejando:
\begin{equation}
  g_n(t) = \frac{2}{R^2\left[J_1(\lambda_n)\right]^2}
  \int_0^R G(t)\, J_0\!\left(\frac{\lambda_n r}{R}\right) r\,\dd r
  \label{eq:gn_formula}
\end{equation}

\medskip
\noindent\textbf{Evaluación de la integral espacial.}
Como $G(t) = -T_0\,\omega\cos(\omega t)$ no depende de $r$, sale de la integral:

\begin{equation}
  g_n(t) = \frac{2\,(-T_0\,\omega\cos(\omega t))}{R^2\left[J_1(\lambda_n)\right]^2}
  \int_0^R J_0\!\left(\frac{\lambda_n r}{R}\right) r\,\dd r
\end{equation}

Usando la \textbf{identidad de integración} del apéndice:
\begin{equation}
  \int_0^R J_0\!\left(\frac{\lambda_n r}{R}\right) r\,\dd r = \frac{R^2}{\lambda_n}\,J_1(\lambda_n)
\end{equation}

se obtiene, cancelando $R^2$ y un factor $J_1(\lambda_n)$:

\begin{resultado}[Coeficiente $g_n(t)$]
\begin{equation}
  g_n(t) = \frac{-2\,T_0\,\omega\cos(\omega t)}{\lambda_n\, J_1(\lambda_n)}
  \label{eq:gn}
\end{equation}
\end{resultado}

\subsection{Parte B: Resolución de la EDO para \texorpdfstring{$T_n(t)$}{Tn(t)}}

Sustituyendo \eqref{eq:gn} en \eqref{eq:ODE_Tn}, la EDO para cada modo es:

\begin{equation}
  T_n'(t) + k_n\, T_n(t) = \frac{-2\,T_0\,\omega}{\lambda_n\,J_1(\lambda_n)}\cos(\omega t)
  \label{eq:ODE_Tn_explicita}
\end{equation}

con condición inicial $T_n(0) = 0$.

\medskip
\noindent\textbf{Solución por factor integrante.} El factor integrante es $\mu(t) = e^{k_n t}$.
Multiplicando \eqref{eq:ODE_Tn_explicita} por $\mu(t)$:

\begin{equation}
  \frac{d}{dt}\!\left[e^{k_n t} T_n(t)\right]
  = \frac{-2\,T_0\,\omega}{\lambda_n\,J_1(\lambda_n)}\,e^{k_n t}\cos(\omega t)
\end{equation}

Integrando ambos lados:
\begin{equation}
  e^{k_n t} T_n(t)
  = \frac{-2\,T_0\,\omega}{\lambda_n\,J_1(\lambda_n)}
    \int e^{k_n t}\cos(\omega t)\,\dd t + C
\end{equation}

\medskip
\noindent\textbf{Integral clave.} Usando la fórmula estándar de integración por partes:
\begin{equation}
  \int e^{at}\cos(\omega t)\,\dd t
  = \frac{e^{at}}{a^2 + \omega^2}\Bigl[a\cos(\omega t) + \omega\sin(\omega t)\Bigr]
\end{equation}

con $a = k_n$:

\begin{equation}
  \int e^{k_n t}\cos(\omega t)\,\dd t
  = \frac{e^{k_n t}}{k_n^2 + \omega^2}
    \Bigl[k_n\cos(\omega t) + \omega\sin(\omega t)\Bigr]
\end{equation}

Sustituyendo y dividiendo por $e^{k_n t}$:

\begin{equation}
  T_n(t) = \frac{-2\,T_0\,\omega}{\lambda_n\,J_1(\lambda_n)(k_n^2 + \omega^2)}
  \Bigl[k_n\cos(\omega t) + \omega\sin(\omega t)\Bigr]
  + C\,e^{-k_n t}
  \label{eq:Tn_general}
\end{equation}

\medskip
\noindent\textbf{Determinación de la constante.} Aplicando $T_n(0) = 0$ en \eqref{eq:Tn_general}:

\begin{equation}
  0 = \frac{-2\,T_0\,\omega\,k_n}{\lambda_n\,J_1(\lambda_n)(k_n^2 + \omega^2)} + C
  \quad\Rightarrow\quad
  C = \frac{2\,T_0\,\omega\,k_n}{\lambda_n\,J_1(\lambda_n)(k_n^2 + \omega^2)}
\end{equation}

\medskip
\noindent Sustituyendo $C$ en \eqref{eq:Tn_general} y factorizando:

\begin{resultado}[Coeficiente temporal $T_n(t)$]
\begin{equation}
  T_n(t) = \frac{2\,T_0\,\omega}{\lambda_n\,J_1(\lambda_n)(k_n^2 + \omega^2)}
  \left[k_n\,e^{-k_n t} - k_n\cos(\omega t) - \omega\sin(\omega t)\right]
  \label{eq:Tn_final}
\end{equation}
\end{resultado}

\subsection{Parte C: Ensamblaje de la solución final}

Recordando que $u(r,t) = v(r,t) + T_0\sin(\omega t)$ y $v = \sum_{n=1}^{\infty} T_n(t)\,\phi_n(r)$,
la solución completa es:

\begin{tcolorbox}[
  enhanced,
  colback=azulclaro!60,
  colframe=azulprimario,
  title={\large\bfseries Solución Completa},
  fonttitle=\bfseries\color{white},
  coltitle=white,
  arc=4mm, boxrule=1pt, breakable
]
\begin{equation}
  \boxed{
  u(r,t) = T_0\sin(\omega t)
  + \sum_{n=1}^{\infty}
  \frac{2\,T_0\,\omega}{\lambda_n\,J_1(\lambda_n)(k_n^2 + \omega^2)}
  \left[k_n\,e^{-k_n t} - k_n\cos(\omega t) - \omega\sin(\omega t)\right]
  J_0\!\left(\frac{\lambda_n\, r}{R}\right)
  }
  \label{eq:solucion_final}
\end{equation}

con $k_n = \alpha^2\!\left(\dfrac{\lambda_n}{R}\right)^{\!2}$ y $\lambda_n$ el $n$-ésimo cero de $J_0$.
\end{tcolorbox}

\subsection{Interpretación física de los términos}

La solución \eqref{eq:solucion_final} tiene una estructura tripartita con significado físico claro:

\begin{table}[h!]
\centering
\renewcommand{\arraystretch}{1.6}
\begin{tabular}{>{\bfseries\color{azulprimario}}p{3.5cm} p{4.5cm} p{5.5cm}}
\toprule
\color{azulprimario}Término & \color{azulprimario}Expresión & \color{azulprimario}Significado físico \\
\midrule
Frontera & $T_0\sin(\omega t)$ & Oscilación impuesta en el borde del disco \\
Régimen permanente & $-[k_n\cos(\omega t) + \omega\sin(\omega t)]$ & Respuesta estacionaria del disco a la excitación armónica \\
Transitorio & $k_n\,e^{-k_n t}$ & Efecto de la condición inicial; decae exponencialmente \\
\bottomrule
\end{tabular}
\caption{Descomposición física de la solución.}
\end{table}

\begin{destacado}
\textbf{Comportamiento asintótico.} A medida que $t \to \infty$, el término transitorio
$e^{-k_n t} \to 0$ para todo $n \geq 1$ (pues $k_n > 0$). El disco \emph{olvida} su
condición inicial y la temperatura converge a la respuesta permanente dictada por la
oscilación de la frontera.
\end{destacado}


% ═══════════════════════════════════════════════════════════════════════════════
\section{Resumen del Método}
% ═══════════════════════════════════════════════════════════════════════════════

El problema fue resuelto mediante el siguiente esquema:

\begin{center}
\begin{tcolorbox}[
  enhanced,
  colback=grisclaro,
  colframe=azulprimario,
  arc=3mm, boxrule=0.6pt,
  width=0.92\textwidth
]
\[
  \underbrace{u\;\text{con CF no homogénea}}_{\text{Problema original}}
  \xrightarrow{\;\text{Paso 1}\;}
  \underbrace{v\;\text{con CF homogénea + }G(t)}_{\text{EDP no homogénea}}
  \xrightarrow{\;\text{Paso 3}\;}
  \underbrace{T_n'(t) + k_n T_n = g_n(t)}_{\text{Familia de EDOs}}
  \xrightarrow{\;\text{Paso 4}\;}
  u(r,t)
\]
\end{tcolorbox}
\end{center}

\begin{enumerate}[leftmargin=2.5em, label=\textbf{Paso \arabic*.}]
  \item \textbf{Transformación:} $v = u - T_0\sin(\omega t)$ homogeneiza la frontera
        e introduce un término fuente $G(t) = -T_0\omega\cos(\omega t)$ en la EDP.

  \item \textbf{Nueva EDP:} La ecuación para $v$ es no homogénea:
        $v_t = \alpha^2\nabla^2 v + G(t)$, con CF y CI nulas.

  \item \textbf{Expansión de Fourier-Bessel:} $v = \sum_n T_n(t)\,J_0(\lambda_n r/R)$
        transforma la EDP en una EDO lineal de primer orden para cada $T_n(t)$.

  \item \textbf{Coeficientes y ensamblaje:} La ortogonalidad determina $g_n(t)$;
        el factor integrante resuelve la EDO; la CI fija la constante de integración.
\end{enumerate}

% ═══════════════════════════════════════════════════════════════════════════════
\appendix
\section{Propiedades de las Funciones de Bessel Utilizadas}
% ═══════════════════════════════════════════════════════════════════════════════

\subsection{Autofunciones y valores propios}

Las soluciones espaciales del problema que satisfacen la condición de nulidad en $r = R$ son:
\[
  \phi_n(r) = J_0\!\left(\frac{\lambda_n r}{R}\right), \quad J_0(\lambda_n) = 0
\]

\subsection{Relación de ortogonalidad}

\begin{equation}
  \int_0^R J_0\!\left(\frac{\lambda_n r}{R}\right) J_0\!\left(\frac{\lambda_m r}{R}\right) r\,\dd r
  =
  \begin{cases}
    0 & n \neq m \\[4pt]
    \dfrac{R^2}{2}\left[J_1(\lambda_n)\right]^2 & n = m
  \end{cases}
\end{equation}

\subsection{Identidad de integración}

\begin{equation}
  \int x\,J_0(x)\,\dd x = x\,J_1(x) + C
\end{equation}

De la cual se deriva, mediante el cambio de variable $x = \lambda_n r/R$:

\begin{equation}
  \int_0^R J_0\!\left(\frac{\lambda_n r}{R}\right) r\,\dd r = \frac{R^2}{\lambda_n}\,J_1(\lambda_n)
\end{equation}

\end{document}
